\section{Conclusion}
\label{sec:conclusion}
This report looked into the application of estimation theory to track the positioning and velocity of a vehicle on a racetrack. The first implementation of the model used the EKF to process Euclidean distances from three beacons around the track. The EKF was able to achieve a execution time within the required sampling period. However it faced some major problems, the reliance of EKF on Gaussian noise distribution means that it is not able to enforce the strict boundaries without contradicting the main assumptions. Montecarlo simulations show how the lateral position frequently leaves the track limits.

To resolve these limitations, a MHE was implemented. Unlike the EKF, the MHE architecture natively supports the strict enforcement of state constraints, such as track boundaries and velocity parameters. This ensures that the estimated trajectory remains strictly confined within the physical limits of the racetrack. Furthermore, by accounting for these constraints within the optimisation window, the MHE elegantly prevents the estimator from becoming mathematically stuck against the track edged a common failure mode observed when applying hard clamps to the EKF.